% Part: sets-functions-relations
% Chapter: infinite
% Section: dedekinds-proof
%
\documentclass[../../../include/open-logic-section]{subfiles}

\begin{document}

\olfileid{sfr}{infinite}{dedekindsproof}

\olsection[ડેડેકિન્ડની ``સાબિતી'']{અનંત ગણના અસ્તિત્વની
ડેડેકિન્ડની ``સાબિતી''}

આ પ્રકરણમાં આપણે ડેડેકિન્ડ બીજગણિતો વડે પ્રાકૃતિક સંખ્યાઓની
ગણસૈદ્ધાંતિક રજૂઆત કરી છે. \olref[arith][ref]{sec}માં આપણે
વિશ્લેષણના અંકગણિતીકરણના (બીજી બાબતો સાથે) તાત્ત્વિક મહત્ત્વ પર
વિચાર કર્યો હતો. હવે આપણે અહીં જે સિદ્ધ કર્યું છે તેના મહત્ત્વ પર
વિચાર કરવો જોઈએ.

સમગ્ર \olref[sfr][arith][]{chap}માં આપણે પ્રાકૃતિક સંખ્યાઓને આપેલી
માની અને તેમનો ઉપયોગ કરીને પૂર્ણાંકો, સંમેય સંખ્યાઓ અને વાસ્તવિક
સંખ્યાઓની સ્પષ્ટ રચના કરી. આ પ્રકરણમાં આપણે પ્રાકૃતિક સંખ્યાઓની સ્પષ્ટ
રચના આપી નથી. આપણે માત્ર એટલું બતાવ્યું છે કે \emph{કોઈ પણ
ડેડેકિન્ડ-અનંત ગણ આપેલો હોય}, તો એવો ગણ વ્યાખ્યાયિત કરી શકાય છે જે
આપણે~$\Nat$ પાસે ઇચ્છીએ છીએ તેવો જ વ્યવહાર કરશે.

તેથી \emph{કઈ વસ્તુઓ પ્રાકૃતિક સંખ્યાઓ છે} એવા અધિભૌતિક પ્રશ્નનો
જવાબ આપ્યો હોવાનો દાવો આપણે કરી શકતા નથી. પરંતુ એ સારી વાત છે. છેવટે,
\olref[sfr][arith][ref]{sec}માં આપણે ભાર મૂક્યો હતો કે
$\Real$ને ડેડેકિન્ડ કાપોના ગણ તરીકે આપેલી વ્યાખ્યાને અનુકૂળ ઠરાવને
બદલે \emph{શોધ} માનવી ખોટી છે. નિર્ણાયક અવલોકન એ છે કે ડેડેકિન્ડ
કાપો વાસ્તવિક સંખ્યાઓના મુખ્ય ગણિતીય ગુણધર્મોને મૂર્ત કરે છે. અહીં
પણ એવું જ છે: \emph{કોઈ પણ} ડેડેકિન્ડ બીજગણિત પ્રાકૃતિક સંખ્યાઓના
મુખ્ય ગણિતીય ગુણધર્મોને મૂર્ત કરે છે. (ખરેખર, બધા ડેડેકિન્ડ બીજગણિતો
\emph{એકરૂપ} છે એવું સાબિત કરીને ડેડેકિન્ડે આ મુદ્દો દૃઢ કર્યો
(\citeyear[Theorems
132--3]{Dedekind1888}). તેથી ઘણા સમકાલીન
``સંરચનાવાદીઓ'' ડેડેકિન્ડને પોતાના પુરોગામી ગણે છે તેમાં નવાઈ નથી.)

 વળી, આપણે બતાવ્યું છે કે પ્રાકૃતિક સંખ્યાઓના સિદ્ધાંતને સહજ સરળ
 ગણસિદ્ધાંતમાં કેવી રીતે સ્થાપિત કરવો. આ ગણસિદ્ધાંત હજી ખાસ્સો
 અનૌપચારિક છે, પરંતુ તેમાં પ્રાકૃતિક સંખ્યાઓ આપેલી છે એવું (દેખીતી
 રીતે) માનવામાં આવતું નથી. તેથી આપણે ડેડેકિન્ડના પોતાના મહત્ત્વાકાંક્ષી
 યોજનાને સાકાર કરવાની દિશામાં હોઈ શકીએ છીએ; તેમણે તેને આ પ્રમાણે
 સમજાવી હતી:
\begin{quote}
	વિજ્ઞાનમાં જે બાબતની સાબિતી શક્ય હોય તે સાબિતી વિના માનવી ન જોઈએ.
	આ માગ વાજબી લાગતી હોવા છતાં, સૌથી સરળ વિજ્ઞાનનો પાયો નાખવાની
	અત્યંત નવીન પદ્ધતિઓમાં પણ તે સંતોષાઈ છે એમ હું માની શકતો નથી;
	અર્થાત્ તર્કશાસ્ત્રનો એ ભાગ જે સંખ્યાઓના સિદ્ધાંત સાથે સંબંધિત છે.
	અંકગણિત (બીજગણિત, વિશ્લેષણ) માત્ર તર્કશાસ્ત્રનો ભાગ છે એમ કહેતાં
	મારો આશય એ છે કે સંખ્યાનો ખ્યાલ અવકાશ અને સમયની કલ્પનાઓ કે
	અંતઃપ્રજ્ઞાથી સંપૂર્ણપણે સ્વતંત્ર છે---હું તેને વિચારના શુદ્ધ
	નિયમોની સીધી નિષ્પત્તિ માનું છું.
	\citep[preface]{Dedekind1888}
\end{quote}
ડેડેકિન્ડનો સાહસિક વિચાર આ છે. માત્ર (સહજ) ગણસિદ્ધાંત વાપરીને
પ્રાકૃતિક સંખ્યાઓ કેવી રીતે રચવી તે આપણે હમણાં જ બતાવ્યું.
\olref[sfr][arith][]{chap}માં પ્રાકૃતિક સંખ્યાઓ અને થોડો ગણસિદ્ધાંત
આપેલાં હોય ત્યારે વાસ્તવિક સંખ્યાઓ કેવી રીતે રચવી તે જોયું. તેથી કદાચ
``અંકગણિત (બીજગણિત, વિશ્લેષણ)'' એ ``માત્ર તર્કશાસ્ત્રનો ભાગ'' ઠરે
(ડેડેકિન્ડના ``તર્કશાસ્ત્ર'' શબ્દના વિસ્તૃત અર્થમાં).

વિચાર આ છે. પરંતુ જરા થોભો. ડેડેકિન્ડ બીજગણિતની આપણી રચના
(પ્રાકૃતિક સંખ્યાઓનો આપણો અવેજ) ડેડેકિન્ડ-અનંત ગણના અસ્તિત્વ પર
આધારિત છે. (ફરી \olref[sfr][infinite][dedekind]{thm:DedekindInfiniteAlgebra}
જુઓ.) જો ``તર્કશાસ્ત્ર'' અથવા ``વિચારના શુદ્ધ નિયમો'' વડે
ડેડેકિન્ડ-અનંત ગણનું અસ્તિત્વ સ્થાપિત ન કરી શકાય, તો યોજના અટકી જાય છે.

તો શું ``વિચારના શુદ્ધ નિયમો'' વડે ડેડેકિન્ડ-અનંત ગણનું અસ્તિત્વ
સ્થાપિત કરી શકાય? ડેડેકિન્ડનો પ્રયાસ આ હતો:
\begin{quote}
  મારા પોતાના વિચારોનું ક્ષેત્ર, એટલે કે મારા વિચારના વિષય બની શકે
  એવી બધી વસ્તુઓની સમગ્રતા $S$, અનંત છે. કારણ કે જો $s$ એ $S$નો
  ઘટક દર્શાવે, તો $s$ મારા વિચારનો વિષય બની શકે છે એવો વિચાર $s'$ પણ
  પોતે $S$નો ઘટક છે. જો આપણે તેને ઘટક~$s$ના પ્રતિબિંબ $\phi(s)$ તરીકે
  લઈએ, તો \dots~$S$ [ડેડેકિન્ડના અર્થમાં] અનંત છે, જે સાબિત કરવાનું હતું.
	\citep[\S66]{Dedekind1888}
\end{quote}
મોટાભાગે અત્યંત ચુસ્ત ગણિતીય સાબિતીઓથી બનેલા પુસ્તકની વચ્ચે આવી વાત
મળે તે ખરેખર આશ્ચર્યજનક છે. બે ટિપ્પણીઓ કરવા જેવી છે.

પહેલી: આ ``સાબિતી''માં આજે આપણે જેને ``ગણિતીય'' સ્વરૂપ કહીએ એવું
ભાગ્યે જ કંઈ છે. તેમાં મનોવૈજ્ઞાનિક વસ્તુઓ (વિચારો)ની વાત થાય છે, અને
તે પણ માત્ર \emph{સંભવિત} વિચારોની.

બીજી: ઓછામાં ઓછું આપણે ડેડેકિન્ડ બીજગણિતો જે રીતે રજૂ કર્યાં છે તે
દૃષ્ટિએ, આ ``સાબિતી''માં સીધી તકનીકી ખામી છે. ડેડેકિન્ડની દલીલ સફળ
હોય તો તે માત્ર એટલું સ્થાપિત કરે છે કે અનંત સંખ્યામાં વસ્તુઓ છે
(ખાસ કરીને, અનંત સંખ્યામાં વિચારો). પરંતુ ડેડેકિન્ડે આપણને એ માનવાનું
કારણ પણ આપવું પડે કે~$S$ એ અનંત ઘણા !!{element}s ધરાવતો એક જ
\emph{ગણ} છે; $S$ને માત્ર \emph{કેટલીક વસ્તુઓ} (બહુવચનમાં) માનવું
યોગ્ય નથી.

ડેડેકિન્ડે અહીં ખામી જોઈ નહીં તે હકીકત કદાચ સૂચવે છે કે
``સમગ્રતા'' શબ્દનો તેમનો ઉપયોગ \emph{ગણ} શબ્દના \emph{આપણા} ઉપયોગને
ચોક્કસપણે અનુસરતો નથી.\footnote{ખરેખર, એવું નહોતું માનવાનાં બીજાં
કારણો પણ છે; \citet[p.~23]{Potter2004} જુઓ.} પરંતુ આમાં બહુ નવાઈ નથી.
છેલ્લાં બે પ્રકરણમાં આપણે હાથ ધરેલો પ્રકલ્પ---પ્રાકૃતિક સંખ્યાઓની
``રચના'', અને તેમાંથી પૂર્ણાંકો, વાસ્તવિક સંખ્યાઓ અને સંમેય સંખ્યાઓની
``રચના''---સંપૂર્ણપણે સહજ રીતે કર્યો છે. કયા ગણો અસ્તિત્વ ધરાવે છે અને
ક્યારે ધરાવે છે તે વિશે ઘણા સામાન્ય સિદ્ધાંતો મૂક્યા વિના, જરૂર પડતી
ગઈ તેમ આ ગણ અથવા તે ગણ માની લીધા છે. પરંતુ આપણે જાણીએ છીએ કે
\emph{કેટલાક} સામાન્ય સિદ્ધાંતો જરૂરી છે; નહિતર આપણે રસેલના
વિરોધાભાસમાં ફસાઈ જઈશું.

 હવે આપણી સહજતાની મર્યાદા છોડવાનો સમય આવી ગયો છે.

\end{document}
